%% ============================================================
%%  SECCIÓN 1 — Introducción
%%  Responsable: Minilux
%% ============================================================

\section{Introducción}

Las moléculas que conforman un líquido se encuentran sometidas a fuerzas de
atracción intermolecular provenientes de sus vecinas. En el interior del fluido,
estas fuerzas se distribuyen simétricamente en todas las direcciones, por lo que su
resultante es nula. Sin embargo, una molécula ubicada en la interfaz
líquido--aire carece de vecinas en la fase gaseosa. En consecuencia, experimenta
una fuerza neta dirigida hacia el interior del líquido y perpendicular a la superficie.
Esta asimetría fundamental dota a la capa superficial de un exceso de energía libre que se
manifiesta macroscópicamente como una tensión: la superficie se comporta como
una membrana elástica que tiende a minimizar su área \cite{serway}.

La magnitud que cuantifica este fenómeno es el \textit{coeficiente de tensión
superficial} $\gamma$, definido como la fuerza por unidad de longitud que ejerce
una porción de superficie sobre la porción adyacente a lo largo de una línea
imaginaria trazada sobre ella:
\begin{equation}
    \gamma = \frac{F}{L},
\end{equation}
donde $F$ es la fuerza tangente a la superficie y $L$ la longitud de la línea de
contacto. Sus unidades en el Sistema Internacional son \unit{\newton\per\meter},
equivalentes a \unit{\joule\per\meter\squared}, lo que pone de manifiesto que
$\gamma$ también representa la energía superficial por unidad de área. Para el
agua pura a \SI{20}{\celsius} el valor aceptado en la bibliografía es
$\gamma_{\text{ref}} = \SI{72.8e-3}{\newton\per\meter}$ \cite{halliday}.

La validez de esta definición se extiende a cualquier interfaz líquida, aunque el
valor de $\gamma$ depende fuertemente de la temperatura y de la presencia de
sustancias tensoactivas. Los surfactantes, como el detergente, poseen moléculas
anfifílicas que se adsorben preferentemente en la superficie y reducen las fuerzas
cohesivas entre las moléculas de agua, disminuyendo $\gamma$ de forma apreciable.
Este efecto tiene relevancia directa en el presente experimento, pues uno de los
métodos empleados utiliza una solución jabonosa \cite{serway}.

El presente experimento determina empíricamente $\gamma$ para el agua mediante dos métodos
independientes basados en el equilibrio estático. En el \textit{primer método} se emplea
una balanza de tipo Mohr--Westphal con un anillo metálico sumergido en agua pura:
la fuerza necesaria para desprender el anillo de la superficie se equilibra con
jinetillos colocados en el brazo mayor, y a partir del torque equivalente se calcula
$\gamma$ usando las circunferencias interna y externa del anillo como longitud de contacto.
En el \textit{segundo método} se forma una película de solución jabonosa sobre un marco de
hilo entre dos tubitos de vidrio; al suspender el sistema, la tensión superficial curva
el hilo inferior y soporta el peso del tubo inferior, lo que permite deducir
$\gamma$ a partir de la geometría y el equilibrio de fuerzas en la película \cite{halliday}.

La comparación entre ambos métodos no solo permite evaluar la precisión instrumental de
cada técnica, sino evidenciar empíricamente la influencia del surfactante
sobre la tensión superficial. Se espera que el primer método, realizado con agua pura,
arroje valores próximos al referencial bibliográfico, mientras que el segundo
registrará valores marcadamente menores por la drástica reducción de las fuerzas
cohesivas inducida por el detergente.